% paper/tables/v2_design.tex -- generated by analyze/paper/tables_paper.py; sources in paper/tables/README.md
\setlength{\tabcolsep}{3pt}\renewcommand{\arraystretch}{0.85}%
\begin{tabular}{@{}p{0.168\linewidth}p{0.46\linewidth}lp{0.25\linewidth}@{}}
\toprule
条件 & 中身 & 組 & 位置づけ \\
\midrule
経験なし・50\%・80\% & A-1 が規則の一歩を選ぶ確率 $p=0$，0.5，0.8 で段階Iを学習 & 組1〜6 & 事前登録（判定は組1〜3） \\
非共有条件 & LSTM を process-1 用と process-2 用に分ける（経験80\%のデータ） & 組1〜3 & 事前登録した対照 \\
符号無作為条件 & 経験80\%と同じ軌跡で，$P_t$ の符号だけをエピソードごとに無作為に反転 & 組1〜3 & 事前登録した対照 \\
零入力条件 & 同じ軌跡で $P_t$ を常に0 & 組1〜3 & 補助の対照（入口の重みは学習されない） \\
好み切替条件 & 好きな色をエピソードごとに赤か緑にし，$P_t$ ＝ 好きな色 $-$ 嫌いな色（規則の確率 0.5） & 組1〜3 & 自分側の $d_P$ は事前登録外の確認（読み方は計算前に記録） \\
写しの比較 & 非共有モデルの process-2 用 LSTM に process-1 用の重みを写す（全部／快・不快の8列，経験あり・なし） & 組1〜3 & 事前登録（G5b）．BA と $\Delta$ は事前登録外 \\
A-2 の規則の確率 0.8 & 段階IIの評価用データで，A-2 が確率 0.8 で好きな色へ向かう & 組1〜6 & 事前登録 \\
\multicolumn{2}{@{}p{0.64\linewidth}}{自分の学習データと同じ作りの評価データ（A-1 は規則の確率 0.8，A-2 は元のまま）} & 組1〜3 & 査読を受けて追加した確認（計算の前に判定の基準を記録） \\
\midrule
指標 & \multicolumn{2}{l}{何を測るか} & 事前登録の扱い \\
\midrule
$g$（G3） & \multicolumn{2}{p{0.535\linewidth}}{入力を隠し快・不快を保つと，予測が赤の方へ進み続ける傾き} & 事前登録した主 \\
$d_P$ & \multicolumn{2}{p{0.535\linewidth}}{$(g_{\mbox{\scriptsize 保つ}}-g_{\mbox{\scriptsize 反転}})/2$．見込みの快・不快への依存} & 事前登録（符号無作為の比較） \\
G5b & \multicolumn{2}{p{0.535\linewidth}}{process-2 の一範囲への $\pm a$ による，予測の中の A-2 のずれの差} & 事前登録した主 \\
符号に従う分 & \multicolumn{2}{p{0.535\linewidth}}{本当の快・不快を保った傾きと，符号を反転した傾きの差の半分} & 事後に定義（組1〜4），組5・6 は事前登録 \\
真値条件の $\Delta$ & \multicolumn{2}{p{0.535\linewidth}}{本当の快・不快と零を入れた，他者側の特徴予測の誤差の差} & 事前登録外（A-2 0.8・符号無作為では事前登録） \\
照合の BA & \multicolumn{2}{p{0.535\linewidth}}{二つの好みの候補から予測誤差の小さい方を選んだ正解率} & 自己視界基準は事前登録した主，他者特徴基準は副 \\
T2・T5 & \multicolumn{2}{p{0.535\linewidth}}{他者用モジュールの出力 $Q2$ が好みや色に応じて変わるか} & 事前登録した主 \\
\bottomrule
\end{tabular}
